\section{Fusion as a Compound Intelligence System}

Fusion turns its design philosophy into a closed loop: understand the state,
identify the missing evidence, allocate the right capability, integrate the
result, and commit one coherent transition. The loop repeats as the task and
the external world change.

\begin{figure}[H]
\centering
\resizebox{0.97\textwidth}{!}{%
\begin{tikzpicture}[
  node distance=0.68cm,
  box/.style={draw=fusionblue!70, rounded corners=3pt, very thick,
    fill=fusionblue!5, minimum height=0.9cm, align=center, inner xsep=9pt},
  source/.style={draw=fusiongray!70, rounded corners=3pt, fill=black!2,
    minimum height=0.9cm, align=center, inner xsep=9pt},
  arrow/.style={-{Latex[length=2.2mm]}, thick, draw=fusiongray}
]
\node[box] (state) {Canonical\\task state};
\node[box, right=of state] (perceive) {Understand\\what is unresolved};
\node[box, right=of perceive] (allocate) {Allocate the\\next capability};
\node[source, right=of allocate] (sources) {Heterogeneous intelligence\\
  Models $\cdot$ Roles $\cdot$ Tools $\cdot$ Verification};
\node[box, right=of sources] (evidence) {Attributed\\evidence};
\node[box, right=of evidence] (commit) {Authoritative\\answer / action};

\draw[arrow] (state) -- (perceive);
\draw[arrow] (perceive) -- (allocate);
\draw[arrow] (allocate) -- (sources);
\draw[arrow] (sources) -- (evidence);
\draw[arrow] (evidence) -- (commit);
\draw[arrow] (commit.south) -- ++(0,-0.42) -| node[pos=0.25,below,
  font=\small,align=center] {new observation or\\committed transition} (state.south);
\end{tikzpicture}
}
\caption{Fusion continuously matches the changing demand of a task with the
available supply of model and tool intelligence.}
\label{fig:system}
\end{figure}

Five system elements make this possible.

\subsection{A canonical task state}

Fusion begins with one normalized history of the task: the user's intent,
constraints, prior decisions, tool observations, produced artifacts, and
acceptance conditions. This state is more than conversation memory. It is the
reference point against which the system determines what has changed, what is
known, and what remains unresolved.

Different components see purpose-built projections of this state. An explorer
may need the open question without the full interaction history. A critic may
need the candidate and its acceptance criteria. The authoritative synthesizer
retains the native conversation and action contract. Bounded views reduce
irrelevant context and accidental authority while making every contribution
auditable against what its source was allowed to observe.

\subsection{Task and progress understanding}

Fusion distinguishes the structure of a task from the progress of a workflow.
Task understanding identifies the kinds of uncertainty present: alternative
reasoning paths, separable coverage regions, domain knowledge, precision
requirements, external observations, and consequential risks. Progress
understanding asks what can happen next: whether the system needs to discover,
execute, verify, answer, or wait for information only the environment can
provide.

This distinction is critical in long workflows. A difficult task does not
justify maximum computation on every turn. Once an artifact exists, the next
valuable operation may be a deterministic test rather than another proposed
solution. Before a consequential action, an independent critique may be worth
more than additional elaboration. Allocation follows the local state of the
work, not a one-time difficulty label.

\subsection{Adaptive capability allocation}

Given the state, Fusion chooses a source, a role, a context view, and a budget.
It also chooses whether to call anything at all. This policy can vary the
number and type of evidence producers, the strength of the authoritative
synthesizer, the depth of independent audit, and the threshold for stopping.

Roles encode the observation the system wants to acquire, not personas for
generating stylistic variety:

\begin{itemize}
  \item \textbf{Solve} develops an independent candidate when the path is uncertain.
  \item \textbf{Explore} searches for discriminating directions when the space is open.
  \item \textbf{Variant} reveals materially different realizations when criteria or preferences are unresolved.
  \item \textbf{Critique} hunts for counterexamples, hidden risks, and failed conditions.
  \item \textbf{Verify} obtains scoped evidence from tools, tests, or checkable artifacts.
\end{itemize}

Because the value of a model is role- and state-dependent, Fusion can use an
asymmetric portfolio deliberately. Efficient models extend the breadth of
search and audit where their contribution is high. Frontier models concentrate
on integration and decisions that genuinely require their capabilities. The
portfolio is not fixed: either class can be admitted or withheld according to
expected value.

\subsection{Evidence, not votes}

Every contribution enters the system as bounded, attributed evidence. A
candidate solution, a criticism, a test result, and a suggested tool call do
not have the same semantics. Fusion preserves these differences rather than
flattening all outputs into interchangeable answers.

The authoritative synthesizer is therefore not a majority voter. It must judge
which claims are load-bearing, reconcile or reject incompatible evidence, and
satisfy the original client contract. Producer suggestions remain advisory;
only the authoritative path can emit the next answer or external action.

\subsection{Continuity across agentic workflows}

For a one-shot request, the loop ends in an answer. For an agentic workflow, it
ends in one state-consistent transition and then begins again. Fusion preserves
the same authority and evidence semantics through four recurring modes:

\begin{equation}
  \textsc{Discover} \rightarrow \textsc{Execute}
  \rightarrow \textsc{Verify} \rightarrow \textsc{Commit}.
\end{equation}

Discovery finds the observation that best separates live hypotheses. Execution
makes one coherent change. Verification connects recorded evidence to explicit
acceptance conditions. Commitment occurs only when material conditions are
supported. This continuity is what allows Fusion to remain effective beyond a
single response: the allocation decision evolves with the task instead of
being reset at every model call.
